% Multi-region transfer experiments (US-073, EPIC 11 Paper Track).
% \input-able section: no \documentclass / \begin{document} here.
% Prose in English (the manuscript is EN). Every number is rendered from the real
% EPIC 12 artefacts by scripts/build_us073_transfer_figures.py; nothing is hand-typed.
% Tables : paper/tables/us-073-transfer/{sen4agrinet_domain_gap,eurocropsml_kshot}.tex
% Figures: paper/figures/us-073-transfer/{kshot_curve,worldcereal_tropical_transfer}.{png,svg}
%          paper/figures/us-075/sen4agrinet_per_class_iou_f1.png
%          paper/figures/us-076/eurocropsml_per_class_f1_vs_k.png
%          figures/us-070/eurocropsml_ae_vs_s2.png
% Required packages in the host preamble: booktabs, graphicx, amsmath, hyperref, url.

\section{Multi-Region Transfer: Train in France, Extend Elsewhere}
\label{sec:multiregion}

A crop-mapping copilot trained only on French labels (PASTIS-R) is, by construction,
a French expert. This section turns that limitation into a measured capability: we
document the \emph{recipe} ``train in France $\rightarrow$ extend elsewhere'' along two
real transfer paths that were executed end-to-end and versioned with DVC, plus an
explicitly out-of-scope tropical
extension. The thesis we defend is deliberately conservative: the model \emph{receives
new classes from other datasets via few-shot adaptation}, and the Franco-Iberian domain
gap is \emph{catastrophic in zero-shot}. Consequently, the \textbf{$\Delta$mIoU between
zero-shot and few-shot is the scientific deliverable}, not a high absolute accuracy. This
is consistent with the limited spatial transferability of frozen Earth-observation
embeddings reported in ``Harvesting AlphaEarth''~\cite{harvesting2026alphaearth}: one must
budget for few-shot finetuning rather than promise zero-shot magic.

All embeddings are AlphaEarth Foundations~\cite{brown2025alphaearth} (Google DeepMind),
accessed through the Google Earth Engine asset \texttt{SATELLITE\_EMBEDDING/V1/ANNUAL},
data version~1.1, 64-dimensional, CC-BY-4.0. The frozen reasoner that turns model outputs
into explanations is Gemini~3.5~Flash~\cite{google2025gemini25} (the deployed cloud default);
following the ``Be My Eyes'' pattern~\cite{huang2025bemyeyes}, our trained models are the
\emph{perceiver} and the LLM is a frozen \emph{reasoner} for communication and data
sovereignty, never a pixel classifier.

Figure~\ref{fig:aoi-italy} situates one of the project's European
areas of interest --- the Italian agricultural region used for the
FarSLIP contrastive distillation (\texttt{farslip-clip-italy-v1}) ---
within the multi-region footprint that spans France (PASTIS-R),
Catalonia (Sen4AgriNet), the EuroCropsML countries
(Estonia/Latvia/Portugal), and the tropical WorldCereal regions
(Brazil, India).

\begin{figure}[htbp]
  \centering
  \includegraphics[width=0.75\linewidth]{figures/us-070/aoi_italy.png}
  \caption{Italian area of interest (Pianura Padana, Po valley) used for
  the FarSLIP contrastive distillation (\texttt{farslip-clip-italy-v1}),
  one node of the multi-region footprint of AgroSatCopilot. Real sampled
  AlphaEarth pixels are coloured by their 64-dimensional annual embedding
  over an OpenStreetMap basemap, to locate the AOI and show that the
  embedding varies spatially across the region. AlphaEarth
  \texttt{SATELLITE\_EMBEDDING/V1/ANNUAL} v1.1
  (CC-BY-4.0)~\cite{brown2025alphaearth}; basemap \textcopyright{}
  OpenStreetMap contributors.}
  \label{fig:aoi-italy}
\end{figure}

\subsection{Label-Space Harmonization (HCAT v3)}
\label{sec:multiregion:labelspace}

Transfer across datasets requires a common label space, otherwise the per-region heads
fight over incompatible gradients. We map the 18 agronomic PASTIS-R classes (the
contiguous \texttt{semantic18} space, ids $0\ldots17$) to HCAT~v3 leaf codes (10-digit,
prefix \texttt{33}~$=$~\texttt{crop\_type}) and collapse them, via the group level of the
HCAT hierarchy, to \textbf{11 canonical macro-classes} (10 crop groups plus a unified
\texttt{void}/background class), well inside the requested 10--20 range. Every HCAT leaf
code is verified to appear verbatim in the EuroCrops HCAT~v3 reference
(\texttt{data/reference/eurocrops\_hcat3.csv}); no code is invented. The mapping and its
collapse rule are documented in \texttt{docs/data/hcat\_crosswalk.md} and materialized as
\texttt{data/reference/hcat\_crosswalk.parquet}.

Three PASTIS classes have no 1:1 HCAT leaf and are flagged \texttt{approx} as explicit
agronomic decisions (not matching failures): \emph{Mixed cereal} maps to the L2 group node
\texttt{cereal}; \emph{Fruits/vegetables/flowers} maps to the dominant representative leaf
\texttt{fresh\_vegetables}; \emph{Leguminous fodder} maps to \texttt{legumes\_harvested\_green}.
The long tail is mitigated by the collapse: \emph{Meadow} ($\sim$45\% of crop parcels) is
isolated as \texttt{grassland}, and the eight sibling cereals (including the
near-inseparable soft and durum winter wheats) are fused into \texttt{cereals}, reducing
the inter-class variance that depressed the 18-class macro-F1. The cross-dataset
convention is uniform: \texttt{background}/\texttt{void} $\rightarrow$
\texttt{ignore\_index = 255} (excluded from the confusion matrix), while a class present in
one dataset but absent in another is treated as \texttt{null-class} (an unknown,
partial-label tag), never as a hard negative.

\subsection{Dense Path: Sen4AgriNet Catalonia (France $\rightarrow$ Spain)}
\label{sec:multiregion:dense}

The best individual French segmenter, TSViT~\cite{tarasiou2023tsvit} with the
phenological branch (TSViT-pheno, fold-4 mIoU~0.6253 / F1~0.7500 on
PASTIS-R~\cite{garnot2021pastis}), is the source model. We finetune it on a small
Sen4AgriNet~\cite{sykas2022sen4agrinet} subset over the Catalonia tile 31TCG
(Sentinel-2, years 2019/2020),
projecting both heads onto the shared HCAT macro space. The subset is intentionally tiny
(10 training patches / 90 tiles; 20 validation patches / 180 tiles), versioned in
DVC/GCS (\texttt{data/sen4agrinet.dvc}); we never download the full multi-terabyte
dataset.

Table~\ref{tab:sen4agrinet_domain_gap} reports the result. The zero-shot transfer yields
an mIoU of exactly~$0.0000$. This is real, not a bug: the French TSViT-pheno projected
$18\rightarrow$macro predicts the Catalonia validation tiles almost exclusively as
\texttt{vegetables} ($\sim$567k px) and \texttt{grassland} ($\sim$416k px), two classes
whose support is \emph{zero} in the Spanish ground truth (dominated instead by
\texttt{cereals}, \texttt{legumes\_fodder}, \texttt{vineyard} and
\texttt{oilseed\_industrial}). With an all-zero confusion diagonal, the mIoU is
exactly~0 -- a textbook catastrophic Franco-Iberian domain gap. We also found that the
Sen4AgriNet Catalonia labels carry S4A crop codes outside the PASTIS-18 nomenclature;
these are honestly routed to \texttt{ignore\_index = 255}, yet the remaining
\texttt{cereals}/\texttt{legumes}/\texttt{vineyard}/\texttt{oilseed} ground truth is
sufficient for a valid measurement. Finetuning the macro head on the 10 Spanish patches
recovers a few-shot mIoU of~$0.2468$ (F1-macro~$0.3005$, pixel-accuracy~$0.9179$), i.e.
$\Delta$mIoU~$= +0.2468$ over 40 epochs. This finetuned model is the
``Franco-Iberian'' segmenter that adds olive, vineyard, sorghum and durum-wheat behaviour
to the French source. The full run is recorded in
\texttt{reports/segmentation/sen4agrinet\_transfer\_result.json} and reproduced by the
executed notebook \texttt{notebooks/segmentation/5c\_transfer\_sen4agrinet.ipynb}.

\input{tables/us-073-transfer/sen4agrinet_domain_gap.tex}

\subsection{Within-Country Transfer: France $\rightarrow$ Brittany (BreizhCrops)}
\label{sec:multiregion:breizhcrops}

Before crossing a border, we ask the cheapest possible transfer question: does the model
hold within the \emph{same country}? Brittany is French but agronomically distinct --- an
oceanic climate, a different crop calendar and a different year --- so it isolates the
region shift from the country shift. We train the project's tabular classifier (XGBoost
\texttt{multi:softprob} over the 185 shared per-parcel features) on PASTIS-R (France, 2019)
and evaluate it \emph{directly}, without retraining, on BreizhCrops~\cite{russwurm2020breizhcrops} (Brittany, 2017; regions \texttt{frh01}$+$\texttt{frh04}),
over the seven crops both datasets share ($64{,}121$ train / $2{,}145$ test parcels).
BreizhCrops carries no AlphaEarth embedding, so this is a raw-feature test, not a
Voting-3 evaluation.

The direct-transfer F1-macro is only $0.2057$ (accuracy $0.2956$): even inside France,
moving region is a real jump. The per-crop split is the honest signal (Table is rendered
from \texttt{data/transfer/pastis\_to\_breizhcrops.parquet}): transfer \emph{holds} where a
crop looks phenologically alike across regions --- rapeseed $0.70$, permanent meadow $0.51$
--- and \emph{collapses} on the cereals whose intra-annual dynamics differ --- maize $0.14$,
barley $0.07$, soft wheat $0.01$, sunflower $0.00$. This within-France result is the
strongest, cheapest illustration of the limited-transferability caveat
(Sec.~\ref{sec:multiregion:caveat}): the failure is driven by the region and its phenology,
not merely by the change of country.

\subsection{Tabular Few-Shot Path: EuroCropsML (Transnational)}
\label{sec:multiregion:fewshot}

The dense path measures a within-Europe domain gap pixel-wise; the tabular path measures
how many local labels close a \emph{transnational} gap. Two factual clarifications are
mandatory. First, \textbf{France is not part of EuroCropsML}: the benchmark~\cite{reuss2025eurocropsml} covers Estonia (EE), Latvia (LV) and Portugal (PT),
so the real protocol is LV$+$PT~$\rightarrow$~EE and LV~$\rightarrow$~EE (Reuss et al.
2025, Table~II), not ``France~$\rightarrow$~Estonia''.
Second, EuroCropsML carries \textbf{no AlphaEarth embedding}: each datapoint is an annual
Sentinel-2 L1C median time series (13 bands, B10 excluded) per parcel. We therefore reuse
the \emph{recipe} of our AlphaEarth tabular baseline (an XGBoost \texttt{multi:softprob}
over a fixed-width per-parcel vector) but build that vector from the parcel's own
Sentinel-2 series; the AlphaEarth-via-GEE variant for these countries stays future work.
Labels are aligned to the same HCAT macro space (the \texttt{EC\_hcat\_c} code is an HCAT
leaf code), yielding 8 macro-classes on this subset, versioned as
\texttt{data/transfer/eurocropsml.dvc}.

Table~\ref{tab:eurocropsml_kshot} and Figure~\ref{fig:eurocropsml_kshot} give the
F1-macro-versus-$k$ curve (mean~$\pm$~std over three seeds) for $k \in
\{1,5,10,20,100,200,500\}$ shots per target class, in three scenarios: source-pre-trained
LV$+$PT~$\rightarrow$~EE, source-pre-trained LV~$\rightarrow$~EE, and a no-pre-train
reference. The reading is the intended one for a domain gap: at $k{=}1$ the no-pre-train
model is near-random ($\sim$0.015 F1-macro) while source pre-training already reaches
$\sim$0.32; as $k$ grows the curves converge, confirming that a handful of local labels --
not a magical zero-shot -- is what closes the gap. The curve is evidence of the measured
transfer cost, never a zero-shot accuracy claim, and is computed strictly over the real
EuroCropsML splits.

\input{tables/us-073-transfer/eurocropsml_kshot.tex}

\begin{figure}[t]
  \centering
  \includegraphics[width=0.9\linewidth]{figures/us-073-transfer/kshot_curve.png}
  \caption{EuroCropsML transnational few-shot curve (LV$[+$PT$]\rightarrow$EE). F1-macro
  on the Estonian query set versus $k$ shots per class, mean over three seeds with std
  error bars. Source pre-training dominates at low $k$ and all curves converge as local
  supervision grows: the gap is closed by few-shot data, not by zero-shot transfer.
  EuroCropsML (Reuss et al. 2025, CC-BY-SA-4.0).}
  \label{fig:eurocropsml_kshot}
\end{figure}

\subsection{Does the Foundation Model Help? AlphaEarth vs.\ Raw Sentinel-2}
\label{sec:multiregion:fm-vs-s2}

The previous curve is built from raw per-parcel Sentinel-2 series because EuroCropsML
ships no embedding. A natural and decisive question is whether the AlphaEarth foundation
model --- the same 64-dimensional annual embedding that underpins the rest of this work
--- actually \emph{closes the transnational domain gap with fewer local labels} than raw
spectral bands. To answer it we re-ran the identical LV$\rightarrow$EE few-shot protocol
(XGBoost \texttt{multi:softprob}, three seeds, nine macro-classes) on two feature spaces:
the parcel's own Sentinel-2 series, and the AlphaEarth zonal embedding pulled for the same
$42{,}473$ EuroCropsML parcels through Google Earth
Engine~(\texttt{SATELLITE\_EMBEDDING/V1/ANNUAL}, v1.1).

Figure~\ref{fig:eurocropsml_ae_vs_s2} reports the head-to-head. \textbf{AlphaEarth wins at
every $k$.} The advantage is largest exactly where it matters --- in the low-label regime:
$\Delta$F1-macro~$=+0.111$ at $k{=}1$ ($0.432$ vs.\ $0.321$) and $+0.087$ at $k{=}5$
($0.436$ vs.\ $0.349$), shrinking to $+0.014$--$0.035$ once hundreds of local shots are
available and both spaces have enough supervision to saturate near $0.53$ at $k{=}2000$.
In other words, \emph{the foundation model closes the same domain gap with roughly an order
of magnitude fewer local labels} than raw bands: a single shot per class on AlphaEarth
already matches what raw Sentinel-2 needs $\sim$20 shots to reach. This is the most
positive transfer result of the section and it is the empirical counterpart to the few-shot
budget warned about in ``Harvesting AlphaEarth''~\cite{harvesting2026alphaearth}: the FM
does not deliver zero-shot magic, but it makes the unavoidable few-shot budget dramatically
cheaper. All numbers are rendered from
\texttt{data/transfer/eurocropsml\_alphaearth\_vs\_s2\_delta.parquet}.

\begin{figure}[t]
  \centering
  \includegraphics[width=0.9\linewidth]{figures/us-070/eurocropsml_ae_vs_s2.png}
  \caption{AlphaEarth foundation-model embedding vs.\ raw Sentinel-2 on the EuroCropsML
  LV$\rightarrow$EE transnational few-shot task (F1-macro vs.\ $k$, mean over three seeds).
  AlphaEarth dominates at every $k$, with the largest margin in the low-label regime
  ($+0.111$ at $k{=}1$, $+0.087$ at $k{=}5$), converging as both spaces saturate near
  $0.53$ at $k{=}2000$: the foundation model closes the domain gap with roughly an order of
  magnitude fewer local labels. Embeddings for the $42{,}473$ parcels pulled via Google
  Earth Engine (\texttt{SATELLITE\_EMBEDDING/V1/ANNUAL} v1.1,
  CC-BY-4.0)~\cite{brown2025alphaearth}; EuroCropsML labels (Reuss et al.\ 2024,
  CC-BY-SA-4.0)~\cite{reuss2025eurocropsml}.}
  \label{fig:eurocropsml_ae_vs_s2}
\end{figure}

\subsection{Per-Class Anatomy: Cardinality Dominates Fragility}
\label{sec:multiregion:perclass}

Macro-averaged scores hide which crops transfer and which collapse. Opening the confusion
matrix per class, on both real paths, exposes a single dominant factor:
\textbf{per-class cardinality, not the abstract domain shift, governs transfer fragility}.
On the EuroCropsML LV$\rightarrow$EE tabular path
(Figure~\ref{fig:eurocropsml_perclass}, left), \texttt{grassland} is already
\emph{saturated} at F1~$\approx0.90$ from $k{=}10$ --- the dominant Baltic land cover has no
gap to close --- and \texttt{cereals} ($0.82\rightarrow0.84$) and
\texttt{oilseed\_industrial} ($0.76\rightarrow0.81$) transfer almost for free. By contrast
\texttt{potato} ($0.35\rightarrow0.50$) and \texttt{legumes\_fodder} ($0.03\rightarrow0.29$)
have a real gap that closes only with hundreds of local labels, while
\texttt{sugar\_beet} ($\sim$16 test parcels), \texttt{orchard} ($\sim$52) and
\texttt{vegetables} ($\sim$66) never clear F1~$0.3$ even at $k{=}500$: their minute support
leaves the feature space too sparse for the $k$-shot to fill. The same signature appears in
the dense Sen4AgriNet path
(Figure~\ref{fig:sen4agrinet_perclass}, right): few-shot \texttt{cereals} recovers
completely (IoU~$0.923$, F1~$0.960$), \texttt{vineyard} shows the textbook
high-precision/low-recall pattern (precision~$0.955$ but recall~$0.151$ --- its signature is
reliable but ten patches do not cover it), and \texttt{oilseed\_industrial} ($2067$~px) and
\texttt{potato} ($164$~px) fail outright (IoU~$0.000$) even after finetuning. We also
measured per-class \emph{negative} transfer: pre-training on neighbouring Latvia helps
shared-phenology crops at $k{=}10$ (\texttt{cereals} $+0.343$, \texttt{oilseed}~$+0.302$,
\texttt{grassland}~$+0.285$) but \emph{hurts} \texttt{legumes\_fodder} ($-0.144$) and
\texttt{vegetables} ($-0.064$), so ``pre-train on the neighbour'' is a per-crop, not a
global, decision. The full tables live in
\texttt{docs/transfer/analisis-clases-transfer.md}.

\begin{figure}[t]
  \centering
  \includegraphics[width=0.49\linewidth]{figures/us-076/eurocropsml_per_class_f1_vs_k.png}
  \hfill
  \includegraphics[width=0.49\linewidth]{figures/us-075/sen4agrinet_per_class_iou_f1.png}
  \caption{Per-class transfer anatomy. \emph{Left:} EuroCropsML LV$\rightarrow$EE F1 per
  macro-class versus $k$ --- robust crops (grassland, cereals, oilseed) saturate with few or
  zero local labels, while low-cardinality crops (sugar\_beet, orchard, vegetables) stay
  below F1~$0.3$ regardless of $k$. \emph{Right:} Sen4AgriNet France$\rightarrow$Catalonia
  zero-shot vs.\ few-shot IoU/F1 per macro-class --- cereals recovers fully, vineyard is
  high-precision/low-recall, and minority classes (oilseed, potato) stay at $0$. In both
  paths the worst classes are the lowest-cardinality ones: support, not the domain shift in
  the abstract, drives fragility. EuroCropsML (CC-BY-SA-4.0)~\cite{reuss2025eurocropsml};
  Sen4AgriNet (CC-BY-SA-4.0)~\cite{sykas2022sen4agrinet}.}
  \label{fig:eurocropsml_perclass}
  \label{fig:sen4agrinet_perclass}
\end{figure}

\subsection{Out of Europe: Tropical Transfer with WorldCereal (Brazil and India)}
\label{sec:multiregion:tropical}

We tested whether the recipe survives a jump to the tropics, where phenology, calendar and
the crop catalogue itself differ from Europe. Using the ESA WorldCereal
reference~\cite{vantricht2023worldcereal} (GEE asset
\texttt{ESA/WorldCereal/2021/MODELS/v100}, CC-BY-4.0) we pulled AlphaEarth embeddings via
GEE for two tropical regions --- the Brazilian Cerrado and Karnataka, India --- and ran two
experiments: a \emph{zero-shot} Europe$\rightarrow$Brazil transfer using the European
classifier directly, and the same few-shot protocol as elsewhere.

The zero-shot result is unambiguous and negative: the European decision boundary
\textbf{does not translate}. Maize, the only class shared with the PASTIS-18 nomenclature,
collapses to F1~$=0.0095$ on Brazil --- effectively zero. The remaining tropical classes
(soybean, generic cropland, non-crop) have no European counterpart at all, so they cannot be
predicted by construction; there is no mapping that turns a French head into a Cerrado head.
This is the strongest tropical statement of the limited-transferability
caveat~\cite{harvesting2026alphaearth}.

Few-shot recovers strongly (Figure~\ref{fig:worldcereal_tropical}). With only $k{=}20$
shots per class, Brazil reaches F1-macro~$=0.626$, about $90\%$ of the in-domain ceiling
($0.724$), and keeps climbing to $0.711$ at $k{=}200$; India is harder but follows the same
shape ($0.495$ at $k{=}20$, $0.656$ at $k{=}200$). The lesson is consistent with the European
paths: a frozen global embedding plus a few dozen local labels per class transfers far better
than a transplanted boundary, but \emph{the local labels are mandatory}. Crucially, the
tropical crop classes do \textbf{not} map onto the European PASTIS-18 leaf space --- only
maize overlaps --- so a truly global copilot needs locally labelled few-shot data in every
climatic zone, not merely more European regions. Numbers are rendered from
\texttt{data/transfer/worldcereal\_fewshot\_results.parquet} and
\texttt{worldcereal\_fewshot\_india.parquet}.

\begin{figure}[t]
  \centering
  \includegraphics[width=0.9\linewidth]{figures/us-073-transfer/worldcereal_tropical_transfer.png}
  \caption{Tropical transfer with WorldCereal AlphaEarth embeddings. Zero-shot
  Europe$\rightarrow$Brazil fails (maize F1~$=0.0095$; the European decision boundary does
  not translate), while few-shot recovers the bulk of the in-domain performance: Brazilian
  Cerrado reaches F1-macro~$0.626$ at $k{=}20$ ($\sim$90\% of the in-domain ceiling $0.724$)
  and $0.711$ at $k{=}200$; Karnataka, India follows the same shape ($0.495$ at $k{=}20$).
  Tropical crop classes do not map onto the European PASTIS-18 leaf space (only maize
  overlaps), so global coverage requires local few-shot supervision per climatic zone, not
  more European regions. AlphaEarth embeddings pulled via Google Earth Engine
  (\texttt{SATELLITE\_EMBEDDING/V1/ANNUAL} v1.1,
  CC-BY-4.0)~\cite{brown2025alphaearth}; reference labels ESA WorldCereal
  (\texttt{ESA/WorldCereal/2021/MODELS/v100}, CC-BY-4.0)~\cite{vantricht2023worldcereal}.}
  \label{fig:worldcereal_tropical}
\end{figure}

\subsection{A Single Multi-Region Classifier: An Honest Negative Result}
\label{sec:multiregion:multiregionmodel}

The preceding paths each adapt to one target. A more ambitious idea --- which we tested and
report here in full, including its failure modes --- is to train \emph{one} classifier over
the AlphaEarth embedding pooling \emph{all} regions at once (PASTIS France, EuroCropsML
Estonia and Latvia, Sen4AgriNet Catalonia and France, WorldCereal Brazil and India). The
hypothesis, due to a team member, is appealing: a crop class that is \emph{weak} in one
region for lack of samples might be \emph{strong} in another that abounds in it, so pooling
should yield a finer, richer taxonomy than any single dataset --- for instance
\texttt{spring\_barley} where PASTIS only says ``cereal''. We present the result honestly:
\textbf{the idea works partially --- it genuinely enlarges the taxonomy --- but it does not
rescue any class above a production-quality threshold, and that negative finding is itself
the contribution.} The full study is in \texttt{docs/transfer/modelo-multiregion.md}.

\paragraph{What works: the taxonomy really grows ($18\rightarrow30$ leaves).}
Harmonising every region to a shared HCAT space, the learnable fine-leaf vocabulary expands
from the $18$ PASTIS leaves to \textbf{30 distinct HCAT leaves}, \emph{with no invented
mappings}: the extra $12$ leaves are contributed by EuroCropsML and are exactly the
distinctions PASTIS collapses --- \texttt{spring\_barley} vs.\ \texttt{winter\_barley},
\texttt{oats}, \texttt{rye}, \texttt{spring}/\texttt{summer\_rapeseed}, \texttt{apples},
\texttt{clover}, \texttt{alfalfa\_lucerne}. Because all regions live in the same
64-dimensional AlphaEarth space, the union is representationally valid even though the label
spaces differ. The fine head is trained only on the $17{,}516$ parcels that carry a real
fine leaf (PASTIS $+$ EuroCropsML); macro-only regions (Sen4AgriNet, WorldCereal) are held
for the collapsed-macro evaluation and never inject a pseudo-coarse class that would
cannibalise the fine leaves.

\paragraph{What does not work: zero classes rescued.}
We deliberately avoid reporting an $N$-class macro-F1 (which falls mechanically as $N$
grows) and instead count how many individual leaf classes cross an F1 threshold, against the
PASTIS-only baseline (Table~\ref{tab:multiregion_threshold}). The multi-region model puts
$7$ leaves over F1~$0.70$ and $2$ over $0.85$; PASTIS-only at the same data budget puts
$6$/$2$, and the full-data PASTIS reference puts $10$/$2$. Defining a \emph{rescued} class as
a leaf that crosses F1~$0.85$ in multi-region and does \emph{not} exist as a learnable class
in the PASTIS-only space, the result is \textbf{0 rescued classes at threshold 0.85}: the two
classes that do cross ($\texttt{winter\_rapeseed\_rape}$, F1~$0.948$; \texttt{vineyard},
F1~$0.905$) were already in PASTIS. The best genuinely new EuroCropsML-only leaf is
\texttt{apples} at F1~$0.686$, then \texttt{spring\_rapeseed} at $0.719$; the
\texttt{spring}/\texttt{winter\_barley} split the model attempts but does not resolve
($0.432$/$0.393$). Adding $12$ fine leaves did not add $12$ good classes.

\begin{table}[t]
  \centering
  \caption{Leaf classes crossing an F1 threshold: the multi-region pooled model vs.\ the
  PASTIS-only baseline. The multi-region taxonomy is wider (30 vs.\ 18 leaves) and does not
  degrade the coarse reading, but it rescues \emph{no} class above the production threshold
  of 0.85 --- the two leaves it places there were already learnable in PASTIS. Numbers from
  \texttt{data/transfer/multiregion\_summary.json}.}
  \label{tab:multiregion_threshold}
  \begin{tabular}{lcccc}
    \toprule
    Model & Total leaves & F1$\ge0.70$ & F1$\ge0.80$ & F1$\ge0.85$ \\
    \midrule
    Multi-region (pooled)            & 30 & \textbf{7} & 2 & 2 \\
    PASTIS-only (capped, same budget) & 18 & 6 & 3 & 2 \\
    PASTIS-only (full, reference)     & 18 & 10 & 4 & 2 \\
    \midrule
    \multicolumn{4}{l}{\emph{Rescued classes ($\ge0.85$, EuroCropsML-exclusive)}} & \textbf{0} \\
    \bottomrule
  \end{tabular}
\end{table}

\paragraph{The real bottleneck: phenological separability, not regions.}
Importantly, the coarse level does \emph{not} degrade: collapsing each fine prediction to its
HCAT macro, the multi-region model scores F1-macro~$=0.658$ on the ten European macro-classes,
on par with the project's mono-region tabular baseline ($0.6535$), despite predicting $30$
leaves instead of $18$. (Including the two tropical macros it cannot predict ---
\texttt{non\_crop} and \texttt{other\_cropland}, both F1~$0.0$ for lack of a leaf that
collapses to them --- the all-class macro-F1 is $0.548$.) The diagnosis is the central
finding: \textbf{the bottleneck is phenological separability, not the number of regions.}
Most new EuroCropsML leaves (oats, rye, clover, alfalfa, spring wheat, fine vegetables) sit at
F1~$0.06$--$0.44$ because they are \emph{intrinsically hard to separate in an annual
embedding}: telling spring barley from winter barley, or oats from wheat, needs the
intra-annual phenology that a single 64-dimensional annual vector averages away. Pooling more
\emph{regions} does not help a class that no region can separate at this temporal resolution;
the indicated next step is multi-date intra-annual temporal features over the cereal/legume
leaves, not more regions. The model remains useful as a richer \emph{macro} vocabulary, but
not yet as a production fine classifier. We report this so the limit is on record, not hidden.

\subsection{Caveat: Limited Spatial Transferability}
\label{sec:multiregion:caveat}

Every real path corroborates the central caveat of ``Harvesting
AlphaEarth''~\cite{harvesting2026alphaearth}: frozen global embeddings transfer poorly
across distant agronomic regions without local adaptation. The dense
France~$\rightarrow$~Catalonia zero-shot mIoU of~$0.0000$ and the tropical
Europe~$\rightarrow$~Brazil zero-shot maize F1 of~$0.0095$ are the strongest possible
empirical statements of this gap, while the EuroCropsML and WorldCereal few-shot curves
quantify the budget needed to bridge it --- and the AlphaEarth-vs-Sentinel-2 comparison
shows the foundation model pays that budget down by roughly an order of magnitude in local
labels ($+0.111$ F1 at $k{=}1$). The honest counterweight is the pooled multi-region
classifier: it widens the taxonomy ($18\rightarrow30$ leaves) without degrading the coarse
reading, yet rescues $0$ classes above F1~$0.85$, because the true bottleneck is
phenological separability in an annual embedding, not the number of regions. We therefore do
\emph{not} claim that zero-shot works outside France, that pooling regions rescues new fine
classes, that the LLM layer ``beats'' Gemini at classification (its role is communication,
explanation and on-premise data sovereignty via Qwen3-30B-A3B~\cite{yang2025qwen3} served
with vLLM GPTQ-Int4 on a single GPU), nor any $\geq$0.80 F1 in the tropics. The
deliverables of this section are the \emph{measured} $\Delta$mIoU, the few-shot curves, the
FM-vs-raw advantage and the negative pooled-model finding --- honest evidence of how, and how
far, the copilot becomes a multi-region expert.

% Share-alike note: figures and reports that combine derivatives of Sen4AgriNet,
% EuroCropsML and WorldCereal inherit the CC-BY(-SA)-4.0 clauses and must be redistributed
% under a compatible license (see docs/licenses/DATA_LICENSE.md). WorldCereal is CC-BY-4.0;
% Sen4AgriNet and EuroCropsML are CC-BY-SA-4.0.
